
The combination consists of three steps.  First, the results are adjusted to a common model to allow 
a consistent comparison of central values and evaluation of uncertainty correlations.  This reference
model includes the description of the $W$-boson production, the Breit-Wigner lineshape, and the $W$-boson
polarization, and is described in Section~\ref{sec:production}.  Second, the correlation of
uncertainties between the experiments is evaluated.  The different center-of-mass energies, initial states,
and lepton pseudorapidity coverage make the correlation non-trivial.
Finally, the results are combined for representative PDF sets, with the compatibility of the measurements
determined for each set.  In addition, other $W$ and $Z$ boson measurements at the Tevatron and LHC are
compared to predictions using these PDF sets, in order to study the reliability of the PDF predictions and
uncertainties for the $m_W$ measurement.

The first step of adjusting each result to a different theoretical model requires an emulation of the 
measurement process, which consists of Monte Carlo event generation (see Sec.~\ref{sec:generators}), detector
simulation (see Sec.~\ref{sec:simulation}), event selection, and a kinematic fit for $m_W$.  The Monte Carlo
samples are produced using a reference value $m_W^{\scriptsize\textrm{ref}}$ for the $W$-boson mass and width
(${\mathrm \Gamma}_W$), and different values of $m_W$ are obtained by reweighting events according to a Breit-Wigner
distribution,
\begin{equation}
  w(m,m_W,m_W^{\scriptsize\textrm{ref}}) = \frac{(m^2 - m_W^2)^2 + m^4 {\mathrm \Gamma}_W^2 / m_W^2} 
{(m^2 - {m_W^{\scriptsize\textrm{ref}}}^2)^2 + m^4 {\mathrm \Gamma}_W^2 /   {m_W^{\scriptsize\textrm{ref}}}^2},
\end{equation}
using the final-state invariant mass $m$.  
This parameterization uses the running width scheme in accordance with 
the published measurement procedures.  The mass reweighting procedure has been checked to give the 
correct target mass value within a statistical uncertainty of $\approx 0.2$~MeV.

The detector simulations used in the original ATLAS, CDF, and D0 measurements are simplified so that
large event samples can be simulated for a variety of PDF sets (see Section~\ref{sec:simulation}).
These simulations do not have the complexity required for a mass measurement in data but are sufficient
for estimating the impact of small theoretical modifications on the measurement.  For the LHCb measurement
no simplification is necessary and the same detector simulation is used as in the original measurement.

The shift in the value of $m_W$ resulting from a change in the generator model is estimated by producing 
template distributions using a given experiment's model, and the same kinematic distributions for an 
alternate model (the “pseudo-data”).  The shift is determined by minimizing the negative log-likelihood
between the pseudo-data and the template distributions.
%This shift is the bias that would
%arise if the alternative model were more correct than the one used in the measurement, so subtracting the shift
%would remove the bias.
In the following we quote the impact $\delta m_W$ of each theoretical shift on a
measurement, i.e. the change in the measured $m_W$ value for a given change in the theoretical model.

A common set of uncertainties and correlations between experiments is obtained by evaluating $\delta m_W$ 
for a variety of PDF sets within the reference theoretical model.  Summing the theoretical shifts gives a
total $\delta m_W$ that we add to each experimental measurement to obtain the value to be used in the
combination.  For each PDF set the combination is performed using the method 
of the best linear unbiased estimator (BLUE)~\cite{Valassi:2003mu}, including both theoretical and experimental 
uncertainties.  The BLUE method is used by the individual experiments to combine results from different 
kinematic distributions, and the combination procedure is validated by reproducing each experiment's published 
value.  Results are presented for a combination of all experimental measurements, as well as for various 
measurement subsets.
